\section{Motivación}
La enseñanza de la programación requiere que el alumno desarrolle habilidades de resolución de problemas lógicos. Con la llegada de los Modelos de Lenguaje Extenso (LLMs), los estudiantes disponen de herramientas capaces de generar código funcional al instante. Sin embargo, esto presenta un riesgo para el aprendizaje profundo, ya que fomenta la copia directa en lugar de la comprensión.

\section{Objetivos}
\subsection{Objetivo General}
El objetivo principal de este trabajo es desarrollar una plataforma web que integre un tutor de programación basado en inteligencia artificial estructurado bajo una metodología socrática.

\subsection{Objetivos Específicos}
\begin{itemize}
    \item Diseñar e implementar el \textit{backend} y \textit{frontend} de una aplicación web interactiva.
    \item Integrar una API de LLM (como OpenAI, Anthropic o local).
    \item Desarrollar el \textit{system prompt} necesario para que la IA actúe exclusivamente como guía y no como un solucionador directo de problemas de código.
    \item Validar la herramienta con usuarios reales para comprobar su eficacia pedagógica.
\end{itemize}

\section{Estructura de la Memoria}
El presente documento se estructura en cinco capítulos principales. Tras esta introducción, el Capítulo~\ref{CAP:ESTADO_ARTE} revisa el estado del arte. Se ha decidido separar el Estado del Arte en su propio capítulo dada la rápida evolución tecnológica de la Inteligencia Artificial Generativa y la necesidad de analizar en profundidad, y de forma aislada, cómo se comparan las distintas herramientas educativas emergentes frente a nuestra propuesta. El Capítulo~\ref{CAP:DISENO} expone la arquitectura del sistema. A continuación, el Capítulo~\ref{CAP:IMPLEMENTACION} detalla el proceso de desarrollo. Finalmente, el Capítulo~\ref{CAP:CONCLUSIONES} expone las conclusiones y líneas de trabajo futuro.
